\documentclass[10pt]{article}
\usepackage[margin=0.85in]{geometry}
\usepackage{amsmath,amssymb,amsthm,hyperref}
\newtheorem{theorem}{Theorem}
\newtheorem{lemma}{Lemma}
\newcommand{\ff}[2]{#1^{\underline{#2}}}
\title{The Gon--Oda determinant conjecture}
\author{Principia Math in Progress}
\date{13 September 2026}
\begin{document}
\maketitle

\paragraph{Task 1. Prove the determinant formulas for \(X_d\).}
Let \(X_d\) be the \((d+1)\)-by-\((d+1)\) coefficient matrix of
Gon and Oda \cite[(7.1)]{GO}, normalized by \(x_{00}=1\).
Put \(t=|m_0|\), \(\varepsilon=\operatorname{sgn}(m_0)i\), and
\(Z_d=(z^{(d)}_{jk})=(\varepsilon^{-j}x_{jk})\).
The source assumes \(t\ge d\) and \(t\equiv d\pmod2\).
We prove its Conjecture 7.3 in this full range.

Here is an explicit characterization of the source matrix that fixes
all its entries without requiring the long sum in (7.1). Set
\[
 P_n(k)=\prod_{\ell=1}^k\frac{2\ell-1}{t-n+2\ell-1}.
\]
Its first two rows, from \cite[(6.19)--(6.20)]{GO}, are
\begin{equation}\label{initial}
 z^{(n)}_{0k}=(-1)^k\binom n{2k}P_n(k),\qquad
 z^{(n)}_{1k}=(-1)^{k+1}\frac tn\binom n{2k-1}P_n(k).
\end{equation}
The remaining rows are uniquely determined by
\begin{equation}\label{rec}
 (j-2k-1)z^{(n)}_{j-1,k}+(t+n-2k+1)z^{(n)}_{j-1,k-1}
 -tz^{(n)}_{jk}+(n-j)z^{(n)}_{j+1,k}=0
\end{equation}
for \(1\le j\le n-1\) and \(0\le k\le n\).
This is exactly \cite[(6.13)]{GO} after dividing by
\(\varepsilon^{j-1}\); \(\varepsilon^2=-1\) fixes the signs.
Indices outside the array are interpreted as zero, and binomial
coefficients outside their usual range are zero. The source also gives
\begin{equation}\label{edges}
 z^{(n)}_{j0}=0\ (j>0),\qquad z^{(n)}_{jn}=0\ (j<n),\qquad
 z^{(n)}_{00}=1.
\end{equation}
In particular, these are the original coefficient matrices, not newly
chosen arrays with the same conjectured determinants.

\begin{theorem}\label{main}
For every natural number \(q\),
\begin{align*}
 \det X_{2q-1}
 &=\varepsilon^{q(2q-1)}
 \frac{\prod_{\ell=1}^{q-1}(t^2-(2\ell-1)^2)^{q-\ell}}
      {\prod_{\ell=1}^{q-1}(t-2\ell)^{2\ell}},\\
 \det X_{2q}
 &=\varepsilon^{q(2q+1)}
 \frac{t^q\prod_{\ell=1}^{q-1}(t^2-(2\ell)^2)^{q-\ell}}
      {\prod_{\ell=1}^{q}(t-2\ell+1)^{2\ell-1}}.
\end{align*}
Consequently \(X_d\) is invertible for every admissible parameter.
\end{theorem}

The proof rests on a reduction of the entire interior matrix.
Let \(K_m\) be the matrix indexed by \(0,\ldots,m\) with
\[
 (K_m)_{r,r-1}=r,\qquad (K_m)_{r,r+1}=m-r,
\]
all other entries being zero. Write \(Z_0=(1)\).

\begin{lemma}\label{reduction}
For every \(d\ge2\),
\begin{equation}\label{dimension}
 (t-d+1)\bigl(z^{(d)}_{jk}\bigr)_{1\le j,k\le d-1}
       =(tI-K_{d-2})Z_{d-2}.
\end{equation}
\end{lemma}

\begin{proof}
The cases \(d=2,3\) follow directly from \eqref{initial}--\eqref{rec}:
their interior matrices are respectively
\(t/(t-1)\) and
\((t-2)^{-1}\left(\begin{smallmatrix}t&-1\\-1&t\end{smallmatrix}\right)\).
Now let \(d\ge4\), put \(m=d-2\), \(h=t-d+1\), and set
\(u_{rc}=z^{(m)}_{rc}\). Define
\begin{equation}\label{vdef}
 v_{r+1,c+1}=h^{-1}\{tu_{rc}-ru_{r-1,c}-(m-r)u_{r+1,c}\},
 \qquad 0\le r,c\le m,
\end{equation}
and set \(v_{j0}=0\). We show that \(v\) equals the interior of \(Z_d\).

First the initial two interior rows agree. The formula for the third
row obtained from \eqref{initial} and \eqref{rec} is
\begin{equation}\label{third}
 z^{(n)}_{2k}=(-1)^k\binom n{2k-2}P_n(k)
 \left(1+\frac{(2k-2)(t^2-n^2)}{(2k-1)n(n-1)}\right),\quad n\ge2.
\end{equation}
We give the algebra for the row comparisons, including vanishing
binomial coefficients. For \(c\ge1\), write
\[
 Q_m(c)=\prod_{\ell=1}^c(t-m+2\ell-1),\qquad
 F_c=\frac{(-1)^c}{2^c c!Q_m(c)},
\]
and use falling factorials, with \(\ff a0=1\). Formula \eqref{initial}
gives
\[
 u_{0c}=F_c\ff m{2c},\qquad
 u_{1c}=-F_c\,2ct\,\ff{(m-1)}{2c-2}.
\]
Hence
\begin{align*}
 tu_{0c}-mu_{1c}
 &=tF_c\{\ff m{2c}+2c\ff m{2c-1}\}\\
 &=tF_c\ff{(m+1)}{2c}
 =h z^{(d)}_{1,c+1}.
\end{align*}
For the next row, \eqref{rec} for \(u\) at \(j=1\) gives
\begin{align*}
 tu_{1c}-u_{0c}-(m-1)u_{2c}
 &=-(2c+1)u_{0c}+(t+m-2c+1)u_{0,c-1}\\
 &=-F_c\ff m{2c-2}\{(m+2)(m-2c+1)+2ct^2\}\\
 &=h z^{(d)}_{2,c+1},
\end{align*}
where the last equality is \eqref{third}. To see the middle
simplification, put \(a=m-2c+1\) and use
\((2c+1)a(a+1)+2c(t^2-a^2)=(m+2)a+2ct^2\).
These are falling-factorial identities, so they remain valid when
some binomial coefficients vanish. At \(c=0\) the same two
equalities read \(t=t\) and \(-1=-1\), by \eqref{edges}.

It remains to propagate these two rows. For any array \(u\), let
\[
 R_m(u;r,c)=(r-2c-1)u_{r-1,c}+(t+m-2c+1)u_{r-1,c-1}
             -tu_{rc}+(m-r)u_{r+1,c}.
\]
Substitution of \eqref{vdef} and collection of coefficients gives the
following identity for arbitrary entries of \(u\):
\begin{align}\label{intertwine}
 hR_{m+2}(v;r+1,c+1)
 ={}&-(r-1)R_m(u;r-1,c)+tR_m(u;r,c)\notag\\
    &-(m-r+1)R_m(u;r+1,c).
\end{align}
For \(1\le r\le m-1\), its right side vanishes by \eqref{rec};
at \(r=1\), the possibly out-of-range equation at index zero has
coefficient zero and is not used. Thus \(v\) satisfies exactly the
row recurrence of \(Z_d\) for \(j=2,\ldots,d-2\) and
\(k=1,\ldots,d-1\). The values at column zero are zero for both
arrays. Since \(d-j\ne0\), equality of the first two interior rows
propagates to every interior row. This proves \eqref{dimension}.
\end{proof}

\begin{proof}[Proof of Theorem \ref{main}]
We first justify the other corner normalization, \(z^{(d)}_{dd}=1\).
It holds directly for \(d=1,2\). For \(d\ge3\), equation \eqref{rec}
at \(j=d-1,k=d\), together with \eqref{edges}, gives
\[
 z^{(d)}_{dd}=-(t-d+1)z^{(d)}_{d-2,d-1}.
\]
The last-column entry of \eqref{dimension} in row \(r=d-3\)
equals \(-z^{(d-2)}_{d-2,d-2}\): the other terms vanish by
\eqref{edges}. Therefore
\(z^{(d)}_{dd}=z^{(d-2)}_{d-2,d-2}=1\), by induction.

Set \(D_d=\det Z_d\). Expansion along the first and last columns,
whose only nonzero entries are the two corner ones, and then
\eqref{dimension}, give
\begin{equation}\label{detrec}
 D_d=\frac{\det(tI-K_{d-2})}{(t-d+1)^{d-1}}D_{d-2},
 \qquad D_0=D_1=1.
\end{equation}
For completeness, the spectrum of \(K_m\) is elementary. Identify
a vector \((a_0,\ldots,a_m)\) with
\(\sum_{r=0}^m\binom mr a_r x^{m-r}y^r\). Then \(K_m\) represents
the operator \(x\partial_y+y\partial_x\). Its eigenbasis
\((x+y)^{m-s}(x-y)^s\), \(0\le s\le m\), has eigenvalues
\(m-2s\). Hence
\[
 \det(tI-K_m)=\prod_{s=0}^m(t-m+2s).
\]
Iterating \eqref{detrec} separately for odd and even \(d\) now gives
\[
 D_{2q-1}=\frac{\prod_{\ell=1}^{q-1}(t^2-(2\ell-1)^2)^{q-\ell}}
                    {\prod_{\ell=1}^{q-1}(t-2\ell)^{2\ell}},\qquad
 D_{2q}=\frac{t^q\prod_{\ell=1}^{q-1}(t^2-(2\ell)^2)^{q-\ell}}
                    {\prod_{\ell=1}^q(t-2\ell+1)^{2\ell-1}}.
\]
Finally, multiplying row \(j\) by \(\varepsilon^j\) gives
\(\det X_d=\varepsilon^{d(d+1)/2}D_d\), with the exponents stated
in the theorem. All numerator factors are positive for \(t\ge d\),
and all denominators are nonzero, proving invertibility.
\end{proof}

\paragraph{Scope.}
The proof covers every dimension, not only the examples \(d\le7\)
computed in the source. It also proves the displayed formulas as
rational identities in \(t\): each dimension has infinitely many
admissible integer values. The symbolic checks supplied with this
manuscript independently compare the reduction with the original
finite-sum entries; they supplement the general proof above.

\paragraph{Task 2. Compute inverse entries.}
Use the interior reduction to construct an explicit entry formula
for \(Z_d^{-1}\), and determine whether it simplifies in terms of
classical orthogonal polynomials.

\paragraph{Task 3. Resolve the singular parameters.}
For each zero of the rational determinant away from its poles,
determine the rank and a basis of the kernel of \(Z_d\). Relate the
nullities to the multiplicities in the determinant product.

\paragraph{Task 4. Bound the conditioning.}
Obtain upper and lower singular-value bounds for \(Z_d(t)\) when
\(t\ge d\), including the regimes \(t=d\) and \(t/d\) fixed as
\(d\) grows.

\paragraph{Task 5. Extend the dimension reduction.}
Investigate corresponding coefficient arrays for principal-series
Siegel--Whittaker functions on \(\mathrm{Sp}(2,\mathbb R)\) and
for the \(\mathrm{SU}(2,2)\) models discussed in the source.
Determine which admit an identity analogous to \eqref{dimension}.

\begin{thebibliography}{1}
\bibitem{GO} Y.~Gon and T.~Oda,
\emph{An explicit integral representation of Siegel--Whittaker
functions on Sp(2,R) for the large discrete series representations},
\href{https://arxiv.org/abs/1412.8330}{arXiv:1412.8330},
Proposition 6.2, equations (6.13), (6.19)--(6.20), and Conjecture 7.3.
The actual 20-page PDF was inspected, especially pp.~13--14 and 18--19.
\end{thebibliography}
\end{document}
